\section{Auto Research abierta: mejora inteligente descentralizada}

\subsection{Conjunto de trabajadores no confiables}

Karpathy está explorando cómo hacer que un \textbf{conjunto de trabajadores no confiables} (untrusted pool of workers) de internet participe en auto research. Idea central:

\begin{importantbox}{La búsqueda es cara, la verificación es barata}
Las mejoras de entrenamiento de LLM tienen una naturaleza similar a blockchain:
\begin{itemize}
    \item Encontrar un buen commit de código puede requerir probar 10,000 ideas (la búsqueda es muy cara)
    \item Pero verificar si un commit es válido solo requiere entrenar el modelo una vez (la verificación es barata)
    \item Un paradigma de cómputo distribuido similar a SETI@Home, Folding@Home
\end{itemize}
"Auto Research at Home": un enjambre de agentes en internet puede colaborar para mejorar los LLM, e incluso es posible que supere a los Frontier Labs.
\end{importantbox}

\subsection{El FLOP como nueva moneda}

Karpathy plantea una reflexión interesante: en el futuro, lo que le importe a la gente quizá no sean los dólares, sino los FLOP. "How many FLOPs do you control instead of what wealth do you control?" Aunque él mismo considera que esto no es del todo cierto, este experimento mental vale la pena considerarlo.

\subsection{Resumen de la sección}

La auto research abierta combina el cómputo distribuido con la investigación en IA, y mediante la naturaleza de "búsqueda cara, verificación barata" permite que incluso un conjunto de trabajadores no confiables participe de forma segura en la mejora de los LLM.

